\section{Space Complexity --- Machines --- Frag Dec}
%%% generated by the blueprint pipeline from TCSlib.Complexity.SpaceComplexity.Machines.FragDec
\subsection*{Space Complexity --- Machines --- Frag Dec: Introduction}

Declarations of the Lean module \texttt{TCSlib.Complexity.SpaceComplexity.Machines.FragDec}.

\subsection*{Space Complexity --- Machines --- Frag Dec: Results}

\begin{definition}[Binary List Decrement]
\lean{Complexity.LogProg.decW}
\label{Complexity.LogProg.decW}
\leanok
Define an operation on finite Boolean lists recursively as follows: it sends the empty
list to itself; on a list beginning with false, it changes that first entry to true and
applies the operation to the remaining list; the one-element list consisting of true is
sent to the empty list; and a list beginning with true and having at least one further
entry has its first entry changed to false, with the rest left unchanged.
\end{definition}

\begin{definition}[Canonical Boolean Lists]
\lean{Complexity.LogProg.Canon}
\label{Complexity.LogProg.Canon}
\leanok
A finite Boolean list is canonical if, whenever it is nonempty, its last entry is true.
\end{definition}

\begin{lemma}[Canonical binary words]
\lean{Complexity.LogProg.canon_bits}
\label{Complexity.LogProg.canon_bits}
\leanok
\uses{Complexity.LogProg.Canon, Complexity.LogProg.bits_getLast}
\difficulty{1}
For every natural number $n$, the bit word produced by its binary expansion is
canonical: if it is nonempty, its last bit is $1$.
\end{lemma}

\begin{lemma}[Canonicality passes to a tail]
\lean{Complexity.LogProg.canon_tail}
\label{Complexity.LogProg.canon_tail}
\leanok
\uses{Complexity.LogProg.Canon}
\difficulty{4}
If a Boolean word formed by prepending one bit $b$ to a word $w$ is canonical, then $w$
is canonical.
\end{lemma}

\begin{lemma}[Decrement after increment]
\lean{Complexity.LogProg.decW_incW}
\label{Complexity.LogProg.decW_incW}
\leanok
\uses{Complexity.LogProg.Canon, Complexity.LogProg.canon_tail, Complexity.LogProg.decW, Complexity.LogProg.incW}
\difficulty{5}
If $w$ is a canonical Boolean word, meaning that every nonempty $w$ ends in $1$, then
decrementing the word obtained by incrementing $w$ returns $w$.
\end{lemma}

\begin{lemma}[Binary expansion of a predecessor]
\lean{Complexity.LogProg.bits_pred}
\label{Complexity.LogProg.bits_pred}
\leanok
\uses{Complexity.LogProg.bits_succ, Complexity.LogProg.canon_bits, Complexity.LogProg.decW, Complexity.LogProg.decW_incW}
\difficulty{4}
For every natural number $n$, the binary bit list of $n-1$ (with subtraction truncated
at zero) is obtained from the binary bit list of $n$ by replacing each initial $0$ by
$1$ until the first $1$, then deleting that $1$; if the list consists only of that $1$,
it becomes empty.
\end{lemma}

\begin{lemma}[Binary expansion after halving]
\lean{Complexity.LogProg.bits_half}
\label{Complexity.LogProg.bits_half}
\leanok
\difficulty{3}
For every natural number $n$, the binary bit list of $\lfloor n/2\rfloor$ is the tail
of the binary bit list of $n$.
\end{lemma}

\begin{lemma}[First bit and parity]
\lean{Complexity.LogProg.bits_head_odd}
\label{Complexity.LogProg.bits_head_odd}
\leanok
\uses{Complexity.LogProg.bits_two_mul_add}
\difficulty{5}
For every natural number $n$, the head of its binary bit list is absent if $n=0$, and
otherwise is the Boolean truth value of the assertion that $n$ is odd.
\end{lemma}

\begin{lemma}[Erasing the last tape cell]
\lean{Complexity.LogProg.bufferTape_erase_last}
\label{Complexity.LogProg.bufferTape_erase_last}
\leanok
\uses{Turing.FinTM.bufferTape}
\difficulty{4}
For every Boolean word $\mathit{pre}$ and bit $b$, changing the tape that stores
$\mathit{pre}$ followed by $b$ so that the cell at position $|\mathit{pre}|$ is blank
yields exactly the tape that stores $\mathit{pre}$.
\end{lemma}

\begin{definition}[Boolean-Dependent Register Action]
\lean{Complexity.LogProg.decDAct}
\label{Complexity.LogProg.decDAct}
\leanok
\uses{Complexity.LogProg.regAct, Turing.Action}
Let $r$ be a work-tape index, let $d_D,d_L,d_B$ be states, and let $rd$ be either a
Boolean symbol or blank. If $rd$ is false, the selected action writes true on tape $r$,
moves that tape's head one cell right, and enters $d_D$. If $rd$ is true, it writes
false on tape $r$, moves that head one cell right, and enters $d_L$. If $rd$ is blank,
it writes nothing on tape $r$, moves that head one cell left, and enters $d_B$.
\end{definition}

\begin{definition}[Blank-Dependent Register Action]
\lean{Complexity.LogProg.decLAct}
\label{Complexity.LogProg.decLAct}
\leanok
\uses{Complexity.LogProg.regAct, Turing.Action}
Let $r$ be a work-tape index, let $d_E,d_B$ be states, and let $rd$ be either a Boolean
symbol or blank. If $rd$ is blank, the selected action writes nothing on tape $r$, moves
that tape's head one cell left, and enters $d_E$; if $rd$ is a Boolean symbol, it writes
nothing on tape $r$, moves that head one cell left, and enters $d_B$.
\end{definition}

\begin{definition}[End action for a register]
\lean{Complexity.LogProg.decEAct}
\label{Complexity.LogProg.decEAct}
\leanok
\uses{Complexity.LogProg.regAct, Turing.Action}
For a register index $r$ among $m$ work tapes and a state $d_B$, let the action write
the blank symbol on tape $r$, move its head one cell to the left, and enter state $d_B$,
leaving the other action components as specified by the register-action construction.
\end{definition}

\begin{lemma}[Run after an initial step]
\lean{Complexity.LogProg.rrun_succ_left}
\label{Complexity.LogProg.rrun_succ_left}
\leanok
\uses{Complexity.LogProg.RProg, Complexity.LogProg.rrun, Complexity.LogProg.rrun_add, Turing.Cfg}
\difficulty{2}
Let $P$ be a register program with oracle, let $c$ be a configuration, and let
$n\in\mathbb N$. The run of length $n+1$ from $c$ equals the run of length $n$ starting
from the configuration reached after one step from $c$.
\end{lemma}

\begin{lemma}[One ordinary transition from a register configuration]
\lean{Complexity.LogProg.rrun_one_reg}
\label{Complexity.LogProg.rrun_one_reg}
\leanok
\uses{Complexity.LogProg.regCfg, Complexity.LogProg.regCfg_state, Complexity.LogProg.rrun, Complexity.LogProg.rrun_one, Complexity.LogProg.rstep_noncall, Turing.Action.apply, Turing.Cfg, Turing.Cfg.inputSymbol, Turing.Cfg.workTapeSymbols, Turing.MultiTapeTM.step}
\difficulty{3}
Let $P$ be a register program with oracle, let $r$ be a register index, and let $c$ be a
configuration. Replace register $r$'s tape by a tape with content $f:\mathbb
Z\to\mathrm{Option}(\mathrm{Bool})$ and place its head at $p\in\mathbb Z$, setting the
state to $s$. If $P$ has no oracle call at $s$, then one step of its run from this
configuration equals applying the machine transition prescribed by $P$ for state $s$,
the scanned input symbol, and the scanned work-tape symbols.
\end{lemma}

\begin{lemma}[Unary decrementing run invariant]
\lean{Complexity.LogProg.decD_run}
\label{Complexity.LogProg.decD_run}
\leanok
\uses{Complexity.LogProg.Canon, Complexity.LogProg.apply_regAct, Complexity.LogProg.bufferTape_erase_last, Complexity.LogProg.canon_tail, Complexity.LogProg.decDAct, Complexity.LogProg.decEAct, Complexity.LogProg.decLAct, Complexity.LogProg.decW, Complexity.LogProg.regCfg, Complexity.LogProg.regCfg_read, Complexity.LogProg.rrun, Complexity.LogProg.rrun_add, Complexity.LogProg.rrun_one_reg, Complexity.LogProg.update_bufferTape_cons, Turing.Cfg, Turing.FinTM.bufferTape}
\difficulty{7}
Write juxtaposition for concatenation of Boolean words. Let $P$ be a register program with an oracle mapping each of $d$ decider indices and each Boolean word to a Boolean answer. Let $r$ be a register and $d_D,d_L,d_E,d_B,\mathrm{next}$ be states. Assume that these four states perform, respectively, the decrement, leftward-decrement, clearing, and increment actions determined by the Boolean symbol read from register $r$, for every input symbol and tuple of work-tape symbols, and that none makes an oracle call. For every initial configuration $c$, canonical Boolean word $w$ (each nonempty canonical word ends in $1$), and Boolean word $\mathrm{pre}$, there are a time $T$ and head position $p$ with $-1\le p<|\mathrm{pre}\,\mathrm{decW}(w)|$ such that the run from state $d_D$, with register $r$ containing $\mathrm{pre}w$ and its head at $|\mathrm{pre}|$, reaches state $d_B$ with register $r$ containing $\mathrm{pre}\,\mathrm{decW}(w)$ and its head at $p$. At every time $t<T$, the run is in state $d_D$, $d_L$, or $d_E$, and its head position $q$ satisfies $|\mathrm{pre}|\le q\le|\mathrm{pre}w|$.
\end{lemma}

\begin{lemma}[Binary decrement computation]
\lean{Complexity.LogProg.dec_run}
\label{Complexity.LogProg.dec_run}
\leanok
\uses{Complexity.LogProg.bits_pred, Complexity.LogProg.canon_bits, Complexity.LogProg.decD_run, Complexity.LogProg.incB_run, Complexity.LogProg.length_bits_mono, Complexity.LogProg.regCfg, Complexity.LogProg.rrun, Complexity.LogProg.rrun_add, Turing.Cfg, Turing.FinTM.bufferTape}
\difficulty{4}
Let $P$ be a register program, let $\mathrm{oracle}$ be a map from decider indices and
Boolean words to Boolean answers, let $r$ be a register, and let
$dD,dL,dE,dB,\mathrm{next}$ be states. Assume the transitions at these states implement,
respectively, the specified binary decrement phases, and assume there are no calls at
these states. For every initial configuration $c$ and natural number $n$, there is a
time $T$ at which the run started in state $dD$ with register $r$ storing the binary
expansion of $n$ and head at zero reaches state $\mathrm{next}$ with that register
storing the binary expansion of $n-1$ (with natural-number truncated subtraction) and
head at zero; at every earlier time it has one of the states $s$, a tape function $f$,
and a head position $q$ such that it is in the corresponding register configuration,
there is no call at $s$, and $-1\le q\le |\operatorname{bits}(n)|$.
\end{lemma}

\begin{definition}[End-of-phase action]
\lean{Complexity.LogProg.toEndAct}
\label{Complexity.LogProg.toEndAct}
\leanok
\uses{Complexity.LogProg.regAct, Turing.Action}
For a register index $r$ among $m$ work tapes, states $c_R$ and $n_x$, and an optional
Boolean value $r_d$, define an action as follows: if $r_d$ is present, leave the work
tape unchanged, move its head one cell to the right, and enter state $c_R$; if $r_d$ is
absent, leave the work tape unchanged, move its head one cell to the left, and enter
state $n_x$. In either case, the other action components are those specified by the
register-action construction.
\end{definition}

\begin{lemma}[Run to the End of a Tape Word]
\lean{Complexity.LogProg.toEnd_run}
\label{Complexity.LogProg.toEnd_run}
\leanok
\uses{Complexity.LogProg.apply_regAct, Complexity.LogProg.regCfg, Complexity.LogProg.regCfg_read, Complexity.LogProg.rrun, Complexity.LogProg.rrun_add, Complexity.LogProg.rrun_one_reg, Complexity.LogProg.toEndAct, Turing.Cfg, Turing.FinTM.bufferTape}
\difficulty{5}
Let $P$ be a register program with oracle, let $r$ be a work-tape register, and let
$c_R$ and $n_x$ be states. Suppose that at state $c_R$ the transition function always
prescribes the specified action that scans register $r$ to the end of its word and
enters state $n_x$, regardless of the scanned input and work-tape symbols, and that
$c_R$ is not a call state. For any initial configuration $c$, word $w$, and integers
$n,p$ with $|w|-p=n$ and $0\le p$, after $n+1$ steps from the configuration with
register $r$ containing $w$ and its head at $p$, the program is in state $n_x$ with that
head at $|w|-1$; at every earlier time $t<n+1$, the configuration still has state $c_R$,
contains the same word on register $r$, and its head is at some $q$ with $p\le q\le
|w|$.
\end{lemma}

\begin{definition}[Clear-end action]
\lean{Complexity.LogProg.clrEAct}
\label{Complexity.LogProg.clrEAct}
\leanok
\uses{Complexity.LogProg.regAct, Turing.Action}
For a register index $r$ among $m$ work tapes, states $c_E$ and $n_{ext}$, and an
optional Boolean value $r_d$, define an action as follows: if $r_d$ is present, write
the blank symbol on tape $r$, move its head one cell to the left, and enter state $c_E$;
if $r_d$ is absent, leave the work tape unchanged, move its head one cell to the right,
and enter state $n_{ext}$. The other action components are those specified by the
register-action construction.
\end{definition}

\begin{lemma}[Register Erasure Run]
\lean{Complexity.LogProg.clrE_run}
\label{Complexity.LogProg.clrE_run}
\leanok
\uses{Complexity.LogProg.apply_regAct, Complexity.LogProg.bufferTape_erase_last, Complexity.LogProg.clrEAct, Complexity.LogProg.regCfg, Complexity.LogProg.regCfg_read, Complexity.LogProg.rrun, Complexity.LogProg.rrun_add, Complexity.LogProg.rrun_one_reg, Turing.Cfg, Turing.FinTM.bufferTape}
\difficulty{6}
Let $P$ be a register program with oracle, choose a register $r$, states
$c_R,c_E,\mathrm{next}$, and suppose that in state $c_R$ every transition is the
prescribed register action, in state $c_E$ every transition is the prescribed erasing
action, and neither state is an oracle-call state. For any configuration $c$ and Boolean
word $w$, for every $k\le |w|$, starting in state $c_E$ with register $r$ containing the
prefix of $w$ of length $k$ and its head at position $k-1$, after $k+1$ steps the
machine is in state $\mathrm{next}$ with register $r$ blank and its head at zero, while
the rest of the configuration is as in $c$. At every earlier time $t<k+1$, the
configuration has state $c_E$ and some tape contents and head position $q$ for register
$r$ satisfying $-1\le q\le k-1$.
\end{lemma}

\begin{lemma}[Register clearing run]
\lean{Complexity.LogProg.clr_run}
\label{Complexity.LogProg.clr_run}
\leanok
\uses{Complexity.LogProg.clrE_run, Complexity.LogProg.regCfg, Complexity.LogProg.rrun, Complexity.LogProg.rrun_add, Complexity.LogProg.toEnd_run, Turing.Cfg, Turing.FinTM.bufferTape}
\difficulty{4}
Let $P$ be a register program with $m$ registers, $d$ deciders, and state set $\Lambda$,
let $\mathcal O$ be an oracle, let $r$ be a register, and let
$c_R,c_E,\mathit{next}\in\Lambda$. Suppose the transition function in state $c_R$ agrees
on every input symbol and tuple of scanned register symbols with the specified action
that begins the clearing run on register $r$, and in state $c_E$ agrees on every such
input and tuple with the specified clearing action leading to $\mathit{next}$. Suppose
also that neither $c_R$ nor $c_E$ is an oracle-call state. For every configuration $c$
and $n\in\mathbb{N}$, there is a time $T$ such that execution from the configuration
obtained from $c$ by setting register $r$ to the tape encoding of the binary digits of
$n$, with its head at $0$ and state $c_R$, reaches the configuration obtained from $c$
by setting that register to the tape encoding of the binary digits of $0$, with its head
at $0$ and state $\mathit{next}$. At each time $t<T$, the configuration has some state
$s$ that is not an oracle-call state, and the head position $q$ of register $r$
satisfies $-1\le q\le$ the length of the binary digit list of $n$.
\end{lemma}
